\section{October 2, 2026}
% Based on pages 4--6 of 6-7-week.pdf.

We derive the leading global error of forward Euler and then construct
its modified equation. These are two descriptions of the same leading
perturbation: one changes the solution, while the other changes the
vector field. We retain the convention that local error is exact
minus numerical, whereas global error is numerical minus exact.

\subsection{Perturbing the starting value of a one-step map}

Consider \(x'=f(t,x)\) and a smooth first-order consistent map
\[
 \Psi^{t+\tau,t}(x)=x+\tau f(t,x)+O(\tau^2).
\]
Assume this expansion also holds after the needed derivatives in
\(x\), uniformly on the region under consideration. Taylor expansion
with respect to the starting value gives
\begin{align}
 \Psi^{t+\tau,t}(x+\delta x)
  &=\Psi^{t+\tau,t}(x)+\delta x
       +\tau f_x(t,x)\delta x\notag\\
  &\quad+O\bigl(\tau^2\|\delta x\|
                      +|\tau|\|\delta x\|^2\bigr).
       \label{eq:step-perturbation-oct-02}
\end{align}
In particular, if \(\delta x=O(\tau^m)\) with integer \(m\geq1\),
the remainder is \(O(\tau^{m+2})\). Here \(f_x=D_xf\) is a
Jacobian, so \(f_x\delta x\) is a matrix--vector product.
The identity term propagates an existing error, and the Jacobian
describes its first-order change over the step.

\subsection{The local error of forward Euler}

Forward Euler is
\begin{equation}\label{eq:euler-oct-02}
 x_{n+1}=x_n+\tau f(t_n,x_n).
\end{equation}
Along an exact solution,
\[
 x''(t)=f_t(t,x(t))+f_x(t,x(t))f(t,x(t)).
\]
Consequently,
\begin{align}
 x(t+\tau)-\Psi^{t+\tau,t}(x(t))
    &=\tau^2 d_2(t)+O(\tau^3),
       \label{eq:euler-defect-oct-02}\\
 d_2(t)&=\frac12
   \bigl(f_t+f_xf\bigr)(t,x(t))=\frac12x''(t).
       \label{eq:euler-d2-oct-02}
\end{align}
The subscript in \(d_2\) records the power of \(\tau\). This is
the leading local-error coefficient denoted by \(d_0\) in the
handwritten notes.

\subsection{An equation for the leading global error}

At grid points, seek an expansion
\begin{equation}\label{eq:euler-global-expansion-oct-02}
 x_\tau(t)=x(t)+\tau C_1(t)+\tau^2C_2(t)+\cdots,
 \qquad x_\tau(t_n)=x_n.
\end{equation}
We use powers beginning at one, so \(C_1\) is the coefficient called
\(c_0\) in the handwritten expansion. Exact initial data give
\(C_j(t_0)=0\) for every correction.

Insert \eqref{eq:euler-global-expansion-oct-02} into the Euler
recurrence. The terms involving \(C_2\) cancel at order \(\tau^2\),
leaving
\[
 \tau^2d_2(t)+\tau\bigl(C_1(t+\tau)-C_1(t)\bigr)
      -\tau^2f_x(t,x(t))C_1(t)=O(\tau^3).
\]
Divide by \(\tau^2\) and compare the leading terms. With
\(A(t)=f_x(t,x(t))\), we obtain
\begin{equation}\label{eq:leading-global-error-oct-02}
 \boxed{C_1'(t)=A(t)C_1(t)-d_2(t),\qquad C_1(t_0)=0.}
\end{equation}
The minus sign agrees with the conventions above: a positive
exact-minus-numerical local defect is subtracted in the recurrence
for the numerical-minus-exact global error.

Let \(W(t,s)\) be the fundamental matrix for the variational
equation,
\[
 \partial_tW(t,s)=A(t)W(t,s),\qquad W(s,s)=I.
\]
Variation of constants gives
\begin{equation}\label{eq:error-integral-oct-02}
 C_1(t)=-\int_{t_0}^t W(t,s)d_2(s)\,ds.
\end{equation}
Thus the leading local defects accumulate after being transported
by the linearized exact flow. If the initial numerical value instead
has expansion \(x_0^{\rm num}=x(t_0)+\tau\eta+O(\tau^2)\), the
additional term is \(W(t,t_0)\eta\).

Higher coefficients satisfy further inhomogeneous variational
equations. Their forcing includes previously computed corrections,
not just the next local-error coefficient. For example, for Euler,
comparison at order \(\tau^3\) gives
\[
 C_2'=A C_2+\frac12f_{xx}[C_1,C_1]
                   -\frac12C_1''-\frac16x''',
 \qquad C_2(t_0)=0,
\]
where all derivatives of \(f\) are evaluated along \(x(t)\).

\begin{example}[Scalar linear equation]
 For \(x'=\lambda x\),
 \(d_2(t)=\lambda^2x(t)/2\) and \(W(t,s)=e^{\lambda(t-s)}\).
 Therefore
 \[
  C_1(t)=-\frac{\lambda^2}{2}(t-t_0)x(t).
 \]
 This agrees with expanding
 \(x_n=(1+\tau\lambda)^n x_0\) at fixed
 \(t_n-t_0=n\tau\).
\end{example}

\subsection{Constructing a modified differential equation}

Backward error analysis seeks a step-dependent vector field
\begin{equation}\label{eq:modified-field-oct-02}
 y'=g_\tau(t,y),\qquad
 g_\tau=f+\tau f_1+\tau^2f_2+\cdots,
 \qquad y(t_0)=x_0,
\end{equation}
whose exact time-\(\tau\) map matches the numerical map order by
order. This is a formal construction. A finite truncation generally
approximates the numerical values at grid points; it does not
interpolate them exactly for all step sizes.

To determine \(f_1\) for Euler, expand one exact step of
\eqref{eq:modified-field-oct-02} from \((t,y)\):
\begin{align*}
 y(t+\tau)
  &=y+\tau g_\tau(t,y)
       +\frac{\tau^2}{2}
          \bigl(\partial_tg_\tau+D_yg_\tau\,g_\tau\bigr)(t,y)
       +O(\tau^3)\\
  &=y+\tau f(t,y)
       +\tau^2\left[f_1(t,y)
                    +\frac12(f_t+f_xf)(t,y)\right]
       +O(\tau^3).
\end{align*}
Euler has no \(\tau^2\) term in its one-step map, so
\begin{equation}\label{eq:euler-modified-oct-02}
 \boxed{f_1=-\frac12(f_t+f_xf),\qquad
 y'=f(t,y)-\frac\tau2(f_t+f_xf)(t,y)+O(\tau^2).}
\end{equation}
The factor \(1/2\) comes from the second Taylor coefficient.
For an autonomous equation this reduces to
\(y'=f(y)-\tau f'(y)f(y)/2+O(\tau^2)\).

The connection with forward error analysis is direct. Expand the
modified solution as \(y=x+\tau y_1+O(\tau^2)\). Then
\[
 y_1'=A(t)y_1+f_1(t,x(t)),\qquad y_1(t_0)=0.
\]
Since \(f_1(t,x(t))=-d_2(t)\), this is precisely
\eqref{eq:leading-global-error-oct-02}; hence \(y_1=C_1\).

If we retain only \(f+\tau f_1\), its exact step and the Euler
step differ by \(O(\tau^3)\). Under uniform smoothness and
stability bounds on a fixed finite interval, the resulting solution
satisfies
\[
 x_n-x(t_n)=O(\tau),\qquad x_n-y(t_n)=O(\tau^2).
\]
The same numerical values are therefore a more accurate approximation
to the solution of this nearby equation.

For \(x'=\lambda x\), the formal modified coefficient can even be
read from the logarithm:
\[
 \frac1\tau\log(1+\tau\lambda)
 =\lambda-\frac\tau2\lambda^2+\frac{\tau^2}{3}\lambda^3-\cdots.
\]
Near \(\tau=0\), its exponential over one step is exactly
\(1+\tau\lambda\), consistent with \eqref{eq:euler-modified-oct-02}.

\subsection{Gradient descent and a modified potential}

Let \(V:\bb R^d\to\bb R\) be smooth and consider the gradient flow
\[
 x'=-\nabla V(x).
\]
The potential decreases along exact trajectories because
\[
 \frac{d}{dt}V(x(t))=-\|\nabla V(x(t))\|_2^2\leq0.
\]
Forward Euler is the gradient descent iteration
\[
 x_{n+1}=x_n-\tau\nabla V(x_n).
\]
Here \(f=-\nabla V\) and \(f'=-\nabla^2V\), so
\[
 f_1=-\frac12\nabla^2V\,\nabla V.
\]
Using \(\nabla\|\nabla V\|_2^2=2\nabla^2V\,\nabla V\), the
first modified equation becomes
\begin{equation}\label{eq:modified-potential-oct-02}
 \boxed{y'=-\nabla V_\tau(y),\qquad
 V_\tau=V+\frac\tau4\|\nabla V\|_2^2.}
\end{equation}
Thus, to first order in the step size, discretization adds a penalty
proportional to the squared gradient. This is the implicit
regularization interpretation mentioned in class. The coefficient is
\(\tau/4\), as required by \eqref{eq:euler-modified-oct-02}.
This describes the truncated modified flow; it does not assert that
Euler decreases either potential for an arbitrary step size.
